\section{Manual (Plataforma web) - Apartado de eventos}

Al ingresar a la página de eventos, se visualizará en primer lugar un listado con todos los eventos activos del club.

\subsection{Reserva de entradas}
El sistema permite reservar una entrada para un evento directamente a un socio. 

Pasos para la reserva:
\begin{itemize}
    \item Seleccionar la opción de reserva de entrada en el evento correspondiente.
    \item Buscar e ingresar el número del socio (por ejemplo, socio número 1150).
    \item Confirmar la operación. El sistema actualizará automáticamente la cantidad de entradas disponibles y el socio pasará a tener la entrada asignada en su perfil.
\end{itemize}

\subsection{Creación de nuevos eventos}
Desde esta misma sección es posible dar de alta nuevos eventos. Este proceso funciona de forma similar a la creación de noticias y requiere completar los siguientes parámetros:

\begin{itemize}
    \item \textbf{Nombre:} Denominación del evento (por ejemplo, evento prueba).
    \item \textbf{Descripción:} Detalles e información adicional del evento.
    \item \textbf{Día y horario:} Fecha del evento, hora de inicio y hora de fin (por ejemplo, 31 de diciembre desde las 20:30 hasta las 22:30).
    \item \textbf{Capacidad máxima:} Límite de asistentes permitidos (por ejemplo, 100 personas).
    \item \textbf{Valor de la entrada:} Precio asignado para asistir, el cual también puede ser cero.
    \item \textbf{Foto:} Imagen ilustrativa del evento (opcional).
\end{itemize}

Una vez creado, el evento aparecerá inmediatamente en el listado principal, quedando habilitado para comenzar a reservar entradas.
